%% Работа с русским языком
% поиск в PDF
\usepackage{cmap}
% русские буквы в формулах
\usepackage{mathtext}
% кодировка
\usepackage[T2A]{fontenc}
% кодировка исходного текста
\usepackage[utf8]{inputenc}
% локализация и переносы
\usepackage[russian]{babel}

%% Пакеты для работы с математикой
\usepackage{amsmath,amsfonts,amssymb,amsthm,mathtools}
\usepackage{icomma}

%% Окружения для утверждений и доказательств
\theoremstyle{plain}
\newtheorem{utverzhdenie}{Утверждение}
\renewcommand{\proofname}{Доказательство}

% %% Номера формул
% % Показывать номера только у тех формул, на которые есть \eqref{} в тексте.
% \mathtoolsset{showonlyrefs=true}
% % Нумерация формул слева
% \usepackage{leqno}

%% Шрифты
% Шрифт "Евклид"
\usepackage{euscript}
% Красивый математический шрифт
\usepackage{mathrsfs}

\usepackage[14pt]{extsizes}

%% Поля (геометрия страницы)
\usepackage[left=3cm,right=1.5cm,top=2cm,bottom=2cm,bindingoffset=0cm]{geometry}

%% Русские списки
\usepackage{enumitem}
\makeatletter
\AddEnumerateCounter{\asbuk}{\russian@alph}{щ}
\makeatother

%% Работа с картинками
\usepackage{caption}
% центрирование подписей к картинкам
\captionsetup{justification=centering}
% для вставки рисунков
\usepackage{graphicx}
% папки с картинками
\graphicspath{{images/}}
% отступ рамки \fbox{} от рисунка
\setlength\fboxsep{3pt}
% толщина линий рамки \fbox{}
\setlength\fboxrule{1pt}
% обтекание рисунков и таблиц текстом
\usepackage{wrapfig}

%% Работа с таблицами
% дополнительная работа с таблицами
\usepackage{array,tabularx,tabulary,booktabs}
% длинные таблицы
\usepackage{longtable}
% слияние строк в таблице
\usepackage{multirow}

%% Красная строка
\setlength{\parindent}{1.25cm}

%% Интервалы
\linespread{1.5}
\usepackage{multirow}

%% TikZ
\usepackage{tikz}
\usetikzlibrary{graphs,graphs.standard}

%% Красивые коробочки
\usepackage{tcolorbox}

%% Перенос знаков в формулах (по Львовскому)
\newcommand*{\hm}[1]{#1\nobreak\discretionary{}{\hbox{$\mathsurround=0pt #1$}}{}}

%% Дополнения
% добавляет возможность работы с командой [H], которая улучшает расположение на странице
\usepackage{float}
% % красивые градусы
% \usepackage{gensymb}
% пакет для подписей к рисункам, в частности, для работы caption*
\usepackage{caption}

% подключаем hyperref (для ссылок внутри  pdf)
\usepackage[unicode, pdftex]{hyperref}

% вставка готового PDF-титульника
\usepackage{pdfpages}

%% Создаём свой неповторимый стиль
\newtcolorbox{tbox}[3][]
{
    colframe = #2!25,
    colback  = #2!10,
    coltitle = #2!20!black,
    title    = {#3},
    #1,
}
